\chapter{Mapping spaces and the exponential law}
\label{chap:mapping}

Spaces of $\Ck{k}$ maps are the models of manifolds of mappings and of the Lie groups of
diffeomorphisms. In this chapter we give them the compact-open $\Ck{k}$ topology and prove the
results that make them usable \cite[\S1.7]{GN26}: superposition operators are continuous and smooth,
pullbacks are continuous linear maps, the differential is a continuous linear map, and the exponential law
$\Ck{k,\ell}(U \times V, F) \cong \Ck{k}(U, \Ck{\ell}(V, F))$ holds, with an embedding in general
and an isomorphism when $V$ is locally compact. Chapter~\ref{chap:limitations} explains why no
better can be expected.

\section{Superposition operators}
\label{sec:pushforward}

\begin{theorem}[Superposition operators are smooth {\cite[Prop.~1.7.1]{GN26}}]
  \label{thm:pushforward}
  \lean{LocallyConvex.Analysis.valuedIn, LocallyConvex.Analysis.isOpen_valuedIn,
    LocallyConvex.Analysis.pushforward, LocallyConvex.Analysis.pushforward_apply,
    LocallyConvex.Analysis.continuousOn_pushforward, LocallyConvex.Analysis.prodCLM,
    LocallyConvex.Analysis.prodCLM_apply, LocallyConvex.Analysis.hasDirDerivAt_pushforward,
    LocallyConvex.Analysis.contDiffMBOn_pushforward,
    LocallyConvex.Analysis.smoothMBOn_pushforward}
  \leanok
  Let $K$ be a compact Hausdorff space, $E$, $F$ locally convex spaces and $g \colon E \to F$
  of class $C^k$ on an open set $U \subseteq E$. With the compact-open topologies, the map
  $C(K, g) \colon C(K, U) \to C(K, F)$, $\gamma \mapsto g \circ \gamma$, is of class $C^k$, with
  $\dd C(K, g)(\gamma, \eta) = (x \mapsto \dd g(\gamma(x), \eta(x)))$.
\end{theorem}

\begin{proof}
  \leanok
  Continuity follows from the continuity of the evaluation $C(K, E) \times K \to E$. For the
  derivative, the difference quotients of $g$ extend continuously to $t = 0$ (Glöckner's lemma),
  and the tube lemma over the compact space $K$ gives convergence in the compact-open topology.
  The derivative is the superposition operator of $\dd g$ on
  $C(K, E \times E) \cong C(K, E) \times C(K, E)$, and the theorem follows by induction on $k$.
\end{proof}

\section{The compact-open \texorpdfstring{$\Ck{k}$}{Ck} topology}
\label{sec:ck-topology}

Let $E$, $F$ be locally convex spaces, $U \subseteq E$ open and $k \in \N$. In Lean, an element
of $\Ck{k}(U, F)$ is a map $E \to F$ that is $\Ck{k}$ on $U$ and vanishes outside $U$; this does
not change the mathematics and makes the family of differentials below point-separating.

\begin{definition}[The compact-open $\Ck{k}$ topology {\cite[Def.~1.7.10]{GN26}}]
  \label{def:ck-topology}
  \lean{LocallyConvex.Analysis.contDiffMBSubmodule, LocallyConvex.Analysis.ContDiffMBMap,
    LocallyConvex.Analysis.ContDiffMBMap.iteratedD,
    LocallyConvex.Analysis.ContDiffMBMap.iteratedDL,
    LocallyConvex.Analysis.ContDiffMBMap.toPi,
    LocallyConvex.Analysis.ContDiffMBMap.instTopologicalSpace,
    LocallyConvex.Analysis.ContDiffMBMap.instIsTopologicalAddGroup,
    LocallyConvex.Analysis.ContDiffMBMap.instContinuousSMulReal,
    LocallyConvex.Analysis.ContDiffMBMap.instLocallyConvexSpaceReal,
    LocallyConvex.Analysis.ContDiffMBMap.instT2Space,
    LocallyConvex.Analysis.ContDiffMBMap.continuous_iff}
  \leanok
  \uses{def:ck, def:iterated, thm:ck-iff-bastiani}
  The compact-open $\Ck{k}$ topology on $\Ck{k}(U, F)$ is the initial topology with respect to
  the linear maps $\dd^{(j)} \colon \Ck{k}(U, F) \to C(U \times E^j, F)$,
  $\gamma \mapsto \dd^{(j)}\gamma$, for $j \leq k$, where $C(U \times E^j, F)$ carries the
  compact-open topology. It makes $\Ck{k}(U, F)$ a locally convex Hausdorff topological vector
  space.
\end{definition}

\begin{lemma}[{\cite[Lemma~1.7.15]{GN26}}, embedding part]
  \label{lem:ck-embedding}
  \lean{LocallyConvex.Analysis.ContDiffMBMap.isEmbedding_toPi}
  \leanok
  \uses{def:ck-topology}
  The map $\Ck{k}(U, F) \to \prod_{j \leq k} C(U \times E^j, F)$,
  $\gamma \mapsto (\dd^{(j)}\gamma)_{j \leq k}$, is a linear topological embedding.
\end{lemma}

\begin{proof}
  \leanok
  This holds by definition of the topology; injectivity comes from $\dd^{(0)}\gamma = \gamma$.
\end{proof}

\begin{lemma}[Limits of differentiable maps]
  \label{lem:limit-differentiable}
  \lean{LocallyConvex.Analysis.hasDirDerivAt_of_tendsto}
  \leanok
  \uses{lem:mean-value-inclusion}
  Let $U \subseteq E$ be open, $y \in E$, and let $(g_i)$ be a net of maps $E \to F$ converging
  pointwise on $U$ to $g$, such that each $g_i$ has directional derivatives $D_i(z)$ in the
  direction $y$ at all $z \in U$, with $D_i$ continuous on $U$. Suppose that $D_i \to D$ in the
  compact-open sense on $U$ (for every compact $K \subseteq U$ and open $O \supseteq D(K)$,
  eventually $D_i(K) \subseteq O$), with $D$ continuous on $U$. Then $g$ has directional
  derivative $D(x)$ in the direction $y$ at every $x \in U$.
\end{lemma}

\begin{proof}
  \leanok
  Let $W$ be a closed convex neighborhood of $D(x)$. For $|s| \leq \delta$, $D(x + sy)$ lies in
  the interior of $W$, hence so does $D_i(x + sy)$ for $i$ large. By the mean value inclusion,
  $t^{-1}(g_i(x + ty) - g_i(x)) \in W$ for $0 < |t| \leq \delta$, and letting $i \to \infty$ the
  same holds for $g$, since $W$ is closed. Glöckner and Neeb use instead their Exercise~1.1.13(b),
  whose hint goes through weak integrals; no completeness is needed here.
\end{proof}

\begin{lemma}[{\cite[Lemma~1.7.15]{GN26}}, closed image]
  \label{lem:ck-closed-image}
  \lean{LocallyConvex.Analysis.ContDiffMBMap.extendPi,
    LocallyConvex.Analysis.ContDiffMBMap.isClosed_range_toPi}
  \leanok
  \uses{lem:ck-embedding, lem:limit-differentiable, thm:ck-iff-bastiani}
  The image of $\Ck{k}(U, F)$ in $\prod_{j \leq k} C(U \times E^j, F)$ is closed.
\end{lemma}

\begin{proof}
  \leanok
  Let $h = (h_j)$ be in the closure of the image and $f := h_0$. By induction on $j$ and
  Lemma~\ref{lem:limit-differentiable}, applied along the filter of $\gamma$ with
  $(\dd^{(j)}\gamma)_j$ close to $h$, the iterated directional derivatives of $f$ exist and
  $\dd^{(j)} f = h_j$. By Bastiani's characterization (Theorem~\ref{thm:ck-iff-bastiani}), $f$
  is $\Ck{k}$.
\end{proof}

\begin{proposition}[Linear operations {\cite[Prop.~1.7.11, Cor.~1.7.13, 1.7.14]{GN26}}]
  \label{prop:ck-linear-operations}
  \lean{LocallyConvex.Analysis.ContDiffMBMap.ofLECLM,
    LocallyConvex.Analysis.ContDiffMBMap.compRightCLM,
    LocallyConvex.Analysis.ContDiffMBMap.compLeftCLM,
    LocallyConvex.Analysis.ContDiffMBMap.prodMkCLM,
    LocallyConvex.Analysis.ContDiffMBMap.prodEquiv}
  \leanok
  \uses{def:ck-topology, thm:chain-rule, lem:prod}
  The following maps are continuous and linear: the inclusion
  $\Ck{l}(U, F) \to \Ck{k}(U, F)$ for $k \leq l$; the pullback $\gamma \mapsto \gamma \circ L$,
  $\Ck{k}(U, F) \to \Ck{k}(L^{-1}(U), F)$, along a continuous linear map $L \colon E_1 \to E$;
  the postcomposition $\gamma \mapsto g \circ \gamma$ with a continuous linear map $g \colon F
  \to G$. Moreover $\Ck{k}(U, F \times G) \cong \Ck{k}(U, F) \times \Ck{k}(U, G)$ as
  topological vector spaces.
\end{proposition}

\begin{proof}
  \leanok
  In each case $\dd^{(j)}$ of the image is obtained from $\dd^{(j)}$ of the argument by pre- or
  postcomposition with a continuous map, or by pairing, and these operations are continuous in
  the compact-open topology (for pairing one restricts to compact subsets of $U \times E^j$,
  which are locally compact).
\end{proof}

\begin{lemma}[The differentials of $\dd\gamma$]
  \label{lem:iterated-d}
  \lean{LocallyConvex.Analysis.ContDiffMBOn.iteratedDMB_dMB}
  \leanok
  \uses{thm:schwarz, thm:iterated-regular, lem:slice}
  Let $\gamma \colon U \to F$ be $\Ck{k+1}$ and $j \leq k$. For $x \in U$ and
  $a, v_1, \dots, v_j, w_1, \dots, w_j \in E$,
  \[
    \dd^{(j)}(\dd\gamma)\bigl((x, a); (v_1, w_1), \dots, (v_j, w_j)\bigr)
      = \dd^{(j+1)}\gamma(x; a, v_1, \dots, v_j)
        + \sum_{i=1}^{j} \dd^{(j)}\gamma(x; v_1, \dots, v_{i-1}, w_i, v_{i+1}, \dots, v_j).
  \]
\end{lemma}

\begin{proof}
  \leanok
  Induction on $j$. In the inductive step, the first summand is differentiated in the direction
  $(v_{j+1}, 0)$ by definition of $\dd^{(j+2)}\gamma$, and in the direction $(0, w_{j+1})$ by
  linearity of $\dd^{(j+1)}\gamma(x; \cdot)$ in its first argument; Schwarz's theorem moves
  $w_{j+1}$ to the last slot. The other summands depend only on $x$.
\end{proof}

\begin{proposition}[The differential is continuous]
  \label{prop:ck-differential}
  \lean{LocallyConvex.Analysis.ContDiffMBMap.dMap, LocallyConvex.Analysis.ContDiffMBMap.dCLM,
    LocallyConvex.Analysis.ContDiffMBMap.iteratedD_dMap,
    LocallyConvex.Analysis.ContDiffMBMap.iteratedD_succ_eq,
    LocallyConvex.Analysis.ContDiffMBMap.continuous_of_continuous_dMap}
  \leanok
  \uses{lem:iterated-d, prop:ck-linear-operations}
  The map $\dd \colon \Ck{k+1}(U, F) \to \Ck{k}(U \times E, F)$, $\gamma \mapsto \dd\gamma$, is
  continuous and linear, and the topology of $\Ck{k+1}(U, F)$ is initial with respect to
  $\gamma \mapsto \gamma \in C(U, F)$ and $\dd$.
\end{proposition}

\begin{proof}
  \leanok
  Continuity of $\dd$ follows from Lemma~\ref{lem:iterated-d}. Conversely,
  $\dd^{(j+1)}\gamma(x; y_0, \dots, y_j) = \dd^{(j)}(\dd\gamma)((x, y_0); (y_1, 0), \dots,
  (y_j, 0))$ (Lemma~\ref{lem:slice}), so each $\dd^{(j+1)}$ is $\dd^{(j)} \circ \dd$ followed by
  a precomposition.
\end{proof}

\begin{theorem}[Superposition operators are continuous {\cite[Prop.~1.7.12, Cor.~1.7.13]{GN26}}]
  \label{thm:ck-superposition}
  \lean{LocallyConvex.Analysis.ContDiffMBMap.compGraph,
    LocallyConvex.Analysis.ContDiffMBMap.dMap_compGraph,
    LocallyConvex.Analysis.ContDiffMBMap.continuous_compGraph,
    LocallyConvex.Analysis.ContDiffMBMap.compLeftOfContDiff,
    LocallyConvex.Analysis.ContDiffMBMap.continuous_compLeftOfContDiff,
    LocallyConvex.Analysis.ContinuousMap.continuous_comp_graph}
  \leanok
  \uses{prop:ck-differential, prop:ck-linear-operations, thm:chain-rule}
  Let $Z$, $E$, $F$ be locally convex spaces, $U \subseteq Z$ open, $r \in \N$ and
  $f \colon Z \times E \to F$ of class $\Ck{r}$ on $U \times E$. Then
  $f_* \colon \Ck{r}(U, E) \to \Ck{r}(U, F)$, $\gamma \mapsto f \circ (\mathrm{id}_U, \gamma)$,
  is continuous. In particular, for $g \colon E \to F$ of class $\Ck{k}$, the map
  $\gamma \mapsto g \circ \gamma$, $\Ck{k}(U, E) \to \Ck{k}(U, F)$, is continuous.
\end{theorem}

\begin{proof}
  \leanok
  Glöckner and Neeb use Faà di Bruno's formula. We argue by induction on $r$ instead. For
  $r = 0$ this is the continuity of $\gamma \mapsto f \circ (\mathrm{id}, \gamma)$ on $C(U, E)$,
  which follows by restricting to compact subsets. By the chain rule,
  $\dd(f_*\gamma)(x, v) = \dd f((x, \gamma(x)), (v, \dd\gamma(x, v)))$, that is,
  $\dd(f_*\gamma) = (\dd f)_*(\gamma \circ \mathrm{pr}_1, \dd\gamma)$ with
  $(\gamma \circ \mathrm{pr}_1, \dd\gamma) \in \Ck{r-1}(U \times Z, E \times E)$ depending
  continuously on $\gamma$ (Propositions~\ref{prop:ck-linear-operations}
  and~\ref{prop:ck-differential}). The induction hypothesis and
  Proposition~\ref{prop:ck-differential} conclude.
\end{proof}

\begin{proposition}[Pullbacks {\cite[Prop.~1.7.11]{GN26}}]
  \label{prop:ck-pullback}
  \lean{LocallyConvex.Analysis.ContDiffMBMap.compRightOfContDiff,
    LocallyConvex.Analysis.ContDiffMBMap.dMap_compRightOfContDiff,
    LocallyConvex.Analysis.ContDiffMBMap.continuous_compRightOfContDiff,
    LocallyConvex.Analysis.ContDiffMBMap.compRightOfContDiffCLM}
  \leanok
  \uses{prop:ck-differential, thm:chain-rule}
  Let $E_1$, $E_2$, $F$ be locally convex spaces, $U \subseteq E_1$ and $V \subseteq E_2$ open,
  and $f \colon E_1 \to E_2$ of class $\Ck{r}$ on $U$ with $f(U) \subseteq V$. Then
  $f^* \colon \Ck{r}(V, F) \to \Ck{r}(U, F)$, $\gamma \mapsto \gamma \circ f$, is continuous and
  linear.
\end{proposition}

\begin{proof}
  \leanok
  Glöckner and Neeb use Faà di Bruno's formula; we argue by induction on $r$, as for
  Theorem~\ref{thm:ck-superposition}. For $r = 0$, $f^*$ is a precomposition. By the chain rule,
  $\dd(\gamma \circ f)(x, v) = \dd\gamma(f(x), \dd f(x, v))$, so $\dd(f^*\gamma) = g^*(\dd\gamma)$
  with $g(x, v) = (f(x), \dd f(x, v))$ of class $\Ck{r-1}$ on $U \times E_1$, and
  Proposition~\ref{prop:ck-differential} and the induction hypothesis conclude.
\end{proof}

\begin{theorem}[The space $\Ck{\infty}(U, F)$ {\cite[Def.~1.7.10, 1.7.11--1.7.15]{GN26}}]
  \label{thm:smooth-space}
  \lean{LocallyConvex.Analysis.smoothMBSubmodule, LocallyConvex.Analysis.SmoothMBMap,
    LocallyConvex.Analysis.SmoothMBMap.toCk, LocallyConvex.Analysis.SmoothMBMap.toPi,
    LocallyConvex.Analysis.SmoothMBMap.ofForall,
    LocallyConvex.Analysis.SmoothMBMap.instLocallyConvexSpaceReal,
    LocallyConvex.Analysis.SmoothMBMap.instT2Space,
    LocallyConvex.Analysis.SmoothMBMap.continuous_iff,
    LocallyConvex.Analysis.SmoothMBMap.isEmbedding_toPi,
    LocallyConvex.Analysis.SmoothMBMap.isClosed_range_toPi,
    LocallyConvex.Analysis.SmoothMBMap.compRightCLM,
    LocallyConvex.Analysis.SmoothMBMap.compLeftCLM,
    LocallyConvex.Analysis.SmoothMBMap.prodEquiv,
    LocallyConvex.Analysis.SmoothMBMap.dCLM,
    LocallyConvex.Analysis.SmoothMBMap.compRightOfContDiffCLM,
    LocallyConvex.Analysis.SmoothMBMap.continuous_compGraph,
    LocallyConvex.Analysis.SmoothMBMap.continuous_compLeftOfContDiff}
  \leanok
  \uses{def:ck-topology, lem:ck-closed-image, prop:ck-linear-operations, prop:ck-differential,
    prop:ck-pullback, thm:ck-superposition}
  The space $\Ck{\infty}(U, F)$ of smooth maps, with the initial topology with respect to all
  the $\dd^{(j)}$, $j \in \N$ (equivalently, with respect to the inclusions into the
  $\Ck{k}(U, F)$), is a locally convex Hausdorff space; it embeds in
  $\prod_{j \in \N} C(U \times E^j, F)$ with closed image, and the results above
  (pullbacks, superposition operators, products, the differential) hold for $k = \infty$.
\end{theorem}

\begin{proof}
  \leanok
  Everything follows from the finite orders, since a map into $\Ck{\infty}(U, F)$ is continuous
  if and only if its compositions with the inclusions into the $\Ck{k}(U, F)$ are. For the
  closed image, the truncations of a point of the closure lie in the closed images of the
  $\Ck{k}(U, F)$, and their preimages all have the same underlying map.
\end{proof}

\section{Smooth superposition operators and locality}
\label{sec:ck-superposition-smooth}

\begin{theorem}[Superposition operators are smooth]
  \label{thm:ck-superposition-smooth}
  \lean{LocallyConvex.Analysis.ContDiffMBMap.compLeftOn,
    LocallyConvex.Analysis.ContDiffMBMap.contDiffMBOn_compLeftOn}
  \leanok
  \uses{thm:ck-superposition, thm:pushforward, lem:ck-closed-image}
  Let $U$ be an open subset of a locally compact locally convex space $Z$, $O \subseteq E$ open and
  $h \colon E \to F$ smooth on $O$. If $g \colon X \to \Ck{r}(U, E)$ is $\Ck{k}$ on an open set
  $W$ and $g(w)(U) \subseteq O$ for $w \in W$, then $w \mapsto h \circ g(w)$ is $\Ck{k}$ on $W$.
\end{theorem}

\begin{proof}
  \leanok
  By induction on $r$. The topology of $\Ck{r}(U, F)$ is initial with respect to the
  differentials $\dd^{(j)}$, and $\Ck{r}(U, F)$ is closed in $\prod_j C(U \times Z^j, F)$, so it
  suffices to treat each $\dd^{(j)}$. For $j = 0$, restriction to the compact subsets of $U$
  reduces the claim to superposition operators on $C(C, E)$ with $C$ compact, which are smooth
  by Theorem~\ref{thm:pushforward}. For $j \ge 1$, $\dd(h \circ \gamma) = \dd h \circ (\gamma \circ
  \mathrm{pr}_1, \dd\gamma)$ is the superposition by the smooth map $\dd h$ of an element of
  $\Ck{r-1}(U \times Z, E \times E)$ that depends continuously and linearly on $\gamma$, and the
  induction hypothesis applies on $U \times Z$.
\end{proof}

\begin{remark}
  Glöckner and Neeb obtain the smoothness of the maps between spaces of $\Ck{\ell}$ maps that they
  need from the exponential law \cite[Proposition~4.1.34]{GN26}. The induction above avoids it. We
  do not claim it is new.
\end{remark}

\begin{proposition}[Locality {\cite[compare Lemma~1.7.43]{GN26}}]
  \label{prop:ck-locality}
  \lean{LocallyConvex.Analysis.ContDiffMBMap.restrictCover,
    LocallyConvex.Analysis.ContDiffMBMap.isEmbedding_restrictCover,
    LocallyConvex.Analysis.ContDiffMBMap.isClosed_range_restrictCover,
    LocallyConvex.Analysis.ContDiffMBMap.contDiffMBOn_of_restrictCLM,
    LocallyConvex.Analysis.ContDiffMBMap.tendsto_of_tendsto_restrictCLM,
    LocallyConvex.Analysis.isEmbedding_restrictCompacts,
    LocallyConvex.Analysis.isClosed_range_restrictCompacts,
    LocallyConvex.Analysis.contDiffMBOn_pi,
    LocallyConvex.Analysis.contDiffMBOn_of_isClosedEmbedding}
  \leanok
  \uses{def:ck-topology}
  \begin{enumerate}
    \item For an open cover $(V_a)$ of $U$, the restrictions
      $\Ck{r}(U, F) \to \prod_a \Ck{r}(V_a, F)$ form a linear topological embedding with closed
      range. Hence being $\Ck{k}$, or converging, in $\Ck{r}(U, F)$ is a local property.
    \item On a locally compact space $Y$, the restrictions $C(Y, F) \to \prod_C C(C, F)$ to the
      compact subsets form a closed embedding.
    \item $\Ck{k}$ maps into a product are the maps with $\Ck{k}$ components, and a map into a
      closed subspace is $\Ck{k}$ as soon as it is $\Ck{k}$ into the ambient space.
  \end{enumerate}
\end{proposition}

\begin{proof}
  \leanok
  (a) The compact-open topology of $C(U \times Z^j, F)$ is initial with respect to the
  restrictions to the $V_a \times Z^j$, since every compact subset of $U$ is covered by finitely
  many compact subsets of the $V_a$; a compatible family glues to a $\Ck{r}$ map. (b) is the case
  of the cover by compact sets, and (c) follows from the definition of the differentials.
\end{proof}

\section{\texorpdfstring{$\Ck{k,\ell}$}{Ckl} maps and the exponential law}
\label{sec:exponential-law}

Let $E_1$, $E_2$, $F$ be locally convex spaces, $U \subseteq E_1$ and $V \subseteq E_2$ open,
and $k, \ell \in \N$.

\begin{definition}[$\Ck{k,\ell}$ maps {\cite[Def.~1.7.17]{GN26}}]
  \label{def:ckl}
  \lean{LocallyConvex.Analysis.iteratedDMB₂, LocallyConvex.Analysis.ContDiffMBOn₂,
    LocallyConvex.Analysis.contDiffMBOn₂_iff}
  \leanok
  \uses{def:iterated, thm:ck-iff-bastiani}
  A map $f \colon U \times V \to F$ is $\Ck{k,\ell}$ if the mixed differentials
  \[
    \dd^{(i,j)} f(x, y; v_1, \dots, v_i; w_1, \dots, w_j)
      = (D_{(v_i,0)} \cdots D_{(v_1,0)} D_{(0,w_j)} \cdots D_{(0,w_1)} f)(x, y)
  \]
  exist for $i \leq k$, $j \leq \ell$ and are continuous on $U \times V \times E_1^i \times E_2^j$.
  Equivalently (by Theorem~\ref{thm:ck-iff-bastiani}), the maps $f(x, \cdot)$ are $\Ck{\ell}$,
  the maps $x \mapsto \dd^{(0,j)} f(x, y; w)$ are $\Ck{k}$, and the $\dd^{(i,j)} f$ are
  continuous.
\end{definition}

\begin{lemma}[Mixed partial derivatives commute, cf.\ {\cite[Lemma~1.3.16]{GN26}}]
  \label{lem:mixed-partials}
  \lean{LocallyConvex.Analysis.eventually_forall_quotient_mem,
    LocallyConvex.Analysis.hasDirDerivAt_comm}
  \leanok
  \uses{lem:limit-differentiable}
  Let $W$ be open and let $g$ have directional derivatives $A$ in the direction $a$ and $B$ in
  the direction $b$ on $W$, with $B$ continuous and with continuous directional derivatives $M$
  in the direction $a$. Then $A$ has directional derivative $M$ in the direction $b$.
\end{lemma}

\begin{proof}
  \leanok
  The difference quotients $r^{-1}(g(z + ra) - g(z))$ converge pointwise to $A$ as $r \to 0$,
  and their derivatives in the direction $b$, $r^{-1}(B(z + ra) - B(z))$, converge to $M$
  uniformly on compact sets. Apply Lemma~\ref{lem:limit-differentiable}. Glöckner and Neeb
  differentiate a weak integral in a completion instead; no completeness is needed here.
\end{proof}

\begin{proposition}[{\cite[Prop.~1.7.21, Cor.~1.7.22]{GN26}}, the case used below]
  \label{prop:ckl-schwarz}
  \lean{LocallyConvex.Analysis.ContDiffMBOn₂.hasDirDerivAt_left,
    LocallyConvex.Analysis.ContDiffMBOn₂.hasDirDerivAt_right,
    LocallyConvex.Analysis.ContDiffMBOn₂.iteratedDMB_right,
    LocallyConvex.Analysis.ContDiffMBOn₂.contDiffMBOn_right_iteratedDMB₂}
  \leanok
  \uses{def:ckl, lem:mixed-partials}
  If $f$ is $\Ck{k,\ell}$, $i \leq k$, $x \in U$ and $v \in E_1^i$, then
  $y \mapsto \dd^{(i,0)} f(x, y; v)$ is $\Ck{\ell}$ on $V$ with $j$-th differential
  $\dd^{(i,j)} f(x, \cdot; v; \cdot)$.
\end{proposition}

\begin{proof}
  \leanok
  By induction on $i$, Lemma~\ref{lem:mixed-partials} moves one derivative in the second
  variable past the derivatives in the first variable.
\end{proof}

\begin{proposition}[Schwarz theorem for $\Ck{k,\ell}$ maps {\cite[Prop.~1.7.21]{GN26}}]
  \label{prop:ckl-schwarz-general}
  \lean{LocallyConvex.Analysis.ContDiffMBOn₂.iteratedDMB_dirOf,
    LocallyConvex.Analysis.ContDiffMBOn₂.hasDirDerivAt_dirOf,
    LocallyConvex.Analysis.ContDiffMBOn₂.iteratedDMB₂_perm_left,
    LocallyConvex.Analysis.ContDiffMBOn₂.iteratedDMB₂_perm_right,
    LocallyConvex.Analysis.ContDiffMBOn₂.iteratedDMB_append_perm}
  \leanok
  \uses{prop:ckl-schwarz}
  Let $f$ be $\Ck{k,\ell}$, $i \leq k$, $j \leq \ell$, $v \in E_1^i$, $w \in E_2^j$ and
  $\pi$ a permutation of $\{1, \ldots, i + j\}$. The iterated directional derivative of $f$ along
  $(v_1, 0), \ldots, (v_i, 0), (0, w_1), \ldots, (0, w_j)$, taken in the order prescribed by
  $\pi$, exists at every point of $U \times V$ and equals $\dd^{(i,j)} f(x, y; v; w)$, which is
  symmetric in $v$ and in $w$.
\end{proposition}

\begin{proof}
  \leanok
  By induction on the number of directions: the last one is either of the form $(v, 0)$, handled
  by the definition of $\dd^{(i,j)} f$, or $(0, w)$, handled by Proposition~\ref{prop:ckl-schwarz}.
  The symmetry in $v$ is the symmetry of the differentials of $x \mapsto \dd^{(0,j)} f(x, y; w)$;
  the symmetry in $w$ follows from Proposition~\ref{prop:ckl-schwarz}.
\end{proof}

\begin{definition}[The space $\Ck{k,\ell}(U \times V, F)$ {\cite[Def.~1.7.17]{GN26}}]
  \label{def:ckl-topology}
  \lean{LocallyConvex.Analysis.ContDiffMBOn₂.add, LocallyConvex.Analysis.ContDiffMBOn₂.const_smul,
    LocallyConvex.Analysis.contDiffMBSubmodule₂, LocallyConvex.Analysis.ContDiffMBMap₂,
    LocallyConvex.Analysis.ContDiffMBMap₂.iteratedD₂,
    LocallyConvex.Analysis.ContDiffMBMap₂.iteratedDL₂, LocallyConvex.Analysis.ContDiffMBMap₂.toPi,
    LocallyConvex.Analysis.ContDiffMBMap₂.instLocallyConvexSpaceReal,
    LocallyConvex.Analysis.ContDiffMBMap₂.instT2Space,
    LocallyConvex.Analysis.ContDiffMBMap₂.isEmbedding_toPi,
    LocallyConvex.Analysis.ContDiffMBMap₂.continuous_iff}
  \leanok
  \uses{def:ckl, def:ck-topology}
  $\Ck{k,\ell}(U \times V, F)$ carries the initial topology with respect to the linear maps
  $\dd^{(i,j)} \colon \Ck{k,\ell}(U \times V, F) \to C(U \times V \times E_1^i \times E_2^j, F)$,
  $i \leq k$, $j \leq \ell$ (compact-open topology on the right). It is a locally convex
  Hausdorff space.
\end{definition}

\begin{proposition}[{\cite[Prop.~1.7.33(a)]{GN26}}]
  \label{prop:ckl-inclusion}
  \lean{LocallyConvex.Analysis.iteratedDMB_slice₂, LocallyConvex.Analysis.iteratedDMB_iteratedDMB,
    LocallyConvex.Analysis.iteratedDMB₂_eq_iteratedDMB,
    LocallyConvex.Analysis.ContDiffMBOn.contDiffMBOn₂, LocallyConvex.Analysis.Opens.prod,
    LocallyConvex.Analysis.ContDiffMBMap₂.ofContDiffMBMap,
    LocallyConvex.Analysis.ContDiffMBMap₂.ofContDiffMBMapCLM}
  \leanok
  \uses{def:ckl-topology, thm:iterated-regular}
  Every $\Ck{k+\ell}$ map on $U \times V$ is $\Ck{k,\ell}$, with
  $\dd^{(i,j)} f(x, y; v; w) = \dd^{(j+i)} f((x, y); (0, w_1), \dots, (0, w_j), (v_1, 0), \dots,
  (v_i, 0))$, and the inclusion $\Ck{k+\ell}(U \times V, F) \to \Ck{k,\ell}(U \times V, F)$ is
  continuous and linear.
\end{proposition}

\begin{proof}
  \leanok
  The formula is an identity between iterated limits; the right hand side is continuous and
  depends continuously on $f$.
\end{proof}

\begin{theorem}[Exponential law {\cite[Thm.~1.7.30]{GN26}}]
  \label{thm:exponential-law}
  \lean{LocallyConvex.Analysis.ContinuousMap.continuous_eval_comp_restrict,
    LocallyConvex.Analysis.ContDiffMBMap₂.partialMap,
    LocallyConvex.Analysis.ContDiffMBMap₂.derivMap,
    LocallyConvex.Analysis.ContDiffMBMap₂.hasDirDerivAt_derivMap,
    LocallyConvex.Analysis.ContDiffMBMap₂.iteratedDMB_partialMap,
    LocallyConvex.Analysis.ContDiffMBMap₂.contDiffMBOn_partialMap,
    LocallyConvex.Analysis.ContDiffMBMap₂.curry,
    LocallyConvex.Analysis.ContDiffMBMap₂.iteratedD_iteratedD_curry,
    LocallyConvex.Analysis.ContDiffMBMap₂.continuous_curry,
    LocallyConvex.Analysis.ContDiffMBMap₂.isEmbedding_curry,
    LocallyConvex.Analysis.ContDiffMBMap₂.curryCLM,
    LocallyConvex.Analysis.ContDiffMBMap₂.uncurry,
    LocallyConvex.Analysis.ContDiffMBMap₂.curryEquiv}
  \leanok
  \uses{prop:ckl-schwarz, def:ckl-topology, lem:mixed-partials, thm:ck-iff-bastiani}
  For $f \in \Ck{k,\ell}(U \times V, F)$, the map $f^\vee \colon U \to \Ck{\ell}(V, F)$,
  $x \mapsto f(x, \cdot)$, is $\Ck{k}$, with
  $\dd^{(j)}(\dd^{(i)} f^\vee(x; v))(y; w) = \dd^{(i,j)} f(x, y; v; w)$, and
  $\Phi \colon \Ck{k,\ell}(U \times V, F) \to \Ck{k}(U, \Ck{\ell}(V, F))$, $f \mapsto f^\vee$, is
  linear and a topological embedding. If $E_2$ is locally compact (for nonempty open $V$, this
  is the same as $V$ being locally compact), $\Phi$ is an isomorphism of topological vector
  spaces, with inverse $g \mapsto ((x, y) \mapsto g(x)(y))$.
\end{theorem}

\begin{proof}
  \leanok
  The candidate for $\dd^{(i)} f^\vee(x; v)$ is $y \mapsto \dd^{(i,0)} f(x, y; v)$, which lies in
  $\Ck{\ell}(V, F)$ by Proposition~\ref{prop:ckl-schwarz}. The difference quotients of
  $x \mapsto \dd^{(i,j)} f(x, \cdot; v; \cdot)$ converge uniformly on compact sets, so the
  derivative exists in $\Ck{\ell}(V, F)$, and Bastiani's characterization shows that $f^\vee$ is
  $\Ck{k}$. Both $\Phi$ and the $\dd^{(i,j)}$ are recovered from each other by evaluations at
  points of compact sets, which are jointly continuous; hence $\Phi$ is an embedding. If $E_2$
  is locally compact, $V \times E_2^j$ is locally compact, so for $g \in \Ck{k}(U, \Ck{\ell}(V, F))$
  the map $(x, y, v, w) \mapsto \dd^{(j)}(\dd^{(i)} g(x; v))(y; w)$ is continuous; it is the
  mixed differential $\dd^{(i,j)}$ of $(x, y) \mapsto g(x)(y)$, which is therefore $\Ck{k,\ell}$.
  Glöckner and Neeb use the evaluation map and a chain rule for $\Ck{r,s}$ maps (with Faà di
  Bruno's formula) instead.
\end{proof}

\begin{proposition}[{\cite[Prop.~1.7.33(b)]{GN26}}]
  \label{prop:ckk-ck}
  \lean{LocallyConvex.Analysis.iteratedDMB_finset_sum, LocallyConvex.Analysis.splitD₁,
    LocallyConvex.Analysis.splitD₂, LocallyConvex.Analysis.splitD,
    LocallyConvex.Analysis.iteratedDMB₂_splitD₁, LocallyConvex.Analysis.iteratedDMB₂_splitD₂,
    LocallyConvex.Analysis.contDiffMBOn₂_splitD,
    LocallyConvex.Analysis.ContDiffMBOn₂.contDiffMBOn_one,
    LocallyConvex.Analysis.ContDiffMBOn₂.contDiffMBOn_prod,
    LocallyConvex.Analysis.ContDiffMBOn₂.congr,
    LocallyConvex.Analysis.ContDiffMBMap₂.splitDMap,
    LocallyConvex.Analysis.ContDiffMBMap₂.continuous_splitDMap,
    LocallyConvex.Analysis.ContDiffMBMap₂.toContDiffMBMap,
    LocallyConvex.Analysis.ContDiffMBMap₂.continuous_toContDiffMBMap,
    LocallyConvex.Analysis.ContDiffMBMap₂.toContDiffMBMapCLM}
  \leanok
  \uses{def:ckl-topology, prop:ckl-schwarz, lem:iterated-d, prop:ck-differential,
    thm:partials}
  Every $\Ck{k,k}$ map on $U \times V$ is $\Ck{k}$, and the inclusion
  $\Ck{k,k}(U \times V, F) \to \Ck{k}(U \times V, F)$ is continuous and linear.
\end{proposition}

\begin{proof}
  \leanok
  Induction on $k$. By the rule on partial differentials, a $\Ck{k+1,k+1}$ map $f$ is $\Ck{1}$
  with $\dd f((x, y), (a, b)) = \dd^{(1,0)} f(x, y; a) + \dd^{(0,1)} f(x, y; b)$. As a function of
  $((x, a), (y, b))$, the right hand side is $\Ck{k,k}$ on $(U \times E_1) \times (V \times E_2)$:
  its mixed differentials are explicit sums of mixed differentials of $f$, by
  Lemma~\ref{lem:iterated-d} and Proposition~\ref{prop:ckl-schwarz}. The induction hypothesis
  gives that $\dd f$ is $\Ck{k}$, and also, with Proposition~\ref{prop:ck-differential}, the
  continuity of the inclusion. Glöckner and Neeb prove instead the explicit formula of
  Proposition~\ref{prop:ckk-formula} for $\dd^{(m)} f$.
\end{proof}

\begin{proposition}[{\cite[Prop.~1.7.33(b), (1.108)]{GN26}}]
  \label{prop:ckk-formula}
  \lean{LocallyConvex.Analysis.ContDiffMBOn₂.iteratedDMB_eq_sum}
  \leanok
  \uses{prop:ckk-ck, prop:ckl-schwarz-general}
  If $f$ is $\Ck{m,m}$, $(x, y) \in U \times V$ and $w \in (E_1 \times E_2)^m$, then
  \[
    \dd^{(m)} f(x, y; w) = \sum_{I \subseteq \{1, \ldots, m\}}
      \dd^{(|I|, m - |I|)} f\bigl(x, y; (w_{i,1})_{i \in I}; (w_{j,2})_{j \notin I}\bigr),
  \]
  with the indices in increasing order.
\end{proposition}

\begin{proof}
  \leanok
  By Proposition~\ref{prop:ckk-ck}, $\dd^{(m)} f(x, y; \cdot)$ is multilinear; expand it along
  $w_a = (w_{a,1}, 0) + (0, w_{a,2})$ and apply Proposition~\ref{prop:ckl-schwarz-general} to each
  term.
\end{proof}

\begin{corollary}[$\Ck{\infty} = \Ck{\infty,\infty}$ {\cite[Cor.~1.7.34]{GN26}}]
  \label{cor:smooth-smooth2}
  \lean{LocallyConvex.Analysis.smoothMBOn_prod_iff, LocallyConvex.Analysis.smoothMBSubmodule₂,
    LocallyConvex.Analysis.SmoothMBMap₂, LocallyConvex.Analysis.SmoothMBMap₂.toCk₂,
    LocallyConvex.Analysis.SmoothMBMap₂.toPi,
    LocallyConvex.Analysis.SmoothMBMap₂.instLocallyConvexSpaceReal,
    LocallyConvex.Analysis.SmoothMBMap₂.continuous_iff,
    LocallyConvex.Analysis.SmoothMBMap₂.equivSmooth}
  \leanok
  \uses{prop:ckk-ck, prop:ckl-inclusion, thm:smooth-space}
  A map on $U \times V$ is smooth if and only if it is $\Ck{k,\ell}$ for all $k, \ell$, and
  $\Ck{\infty}(U \times V, F) = \Ck{\infty,\infty}(U \times V, F)$ as locally convex spaces.
\end{corollary}

\begin{proof}
  \leanok
  Combine Propositions~\ref{prop:ckl-inclusion} and~\ref{prop:ckk-ck}.
\end{proof}

\begin{lemma}[$\Ck{k}$ maps into $\Ck{\infty}(V, F)$]
  \label{lem:ck-into-smooth}
  \lean{LocallyConvex.Analysis.SmoothMBMap.continuousOn_iff,
    LocallyConvex.Analysis.SmoothMBMap.contDiffMBOn_iff,
    LocallyConvex.Analysis.SmoothMBMap.smoothMBOn_iff}
  \leanok
  \uses{thm:smooth-space}
  A map $g \colon W \to \Ck{\infty}(V, F)$ is $\Ck{k}$ if and only if its compositions with the
  inclusions $\Ck{\infty}(V, F) \to \Ck{\ell}(V, F)$ are $\Ck{k}$, for all $\ell$.
\end{lemma}

\begin{proof}
  \leanok
  Induction on $k$. The derivatives of the compositions have the same underlying maps, hence
  define a candidate in $\Ck{\infty}(V, F)$, and the topology of $\Ck{\infty}(V, F)$ is initial
  with respect to the inclusions.
\end{proof}

\begin{theorem}[Exponential law for smooth maps {\cite[Cor.~1.7.35, Rem.~1.7.31(b)]{GN26}}]
  \label{thm:smooth-exponential-law}
  \lean{LocallyConvex.Analysis.SmoothMBMap.toCk₂, LocallyConvex.Analysis.SmoothMBMap.partialMap,
    LocallyConvex.Analysis.SmoothMBMap.smoothMBOn_partialMap,
    LocallyConvex.Analysis.SmoothMBMap.curry,
    LocallyConvex.Analysis.SmoothMBMap.iteratedD_iteratedD_curry,
    LocallyConvex.Analysis.SmoothMBMap.continuous_curry,
    LocallyConvex.Analysis.SmoothMBMap.isEmbedding_curry,
    LocallyConvex.Analysis.SmoothMBMap.curryCLM, LocallyConvex.Analysis.SmoothMBMap.uncurry,
    LocallyConvex.Analysis.SmoothMBMap.curryEquiv,
    LocallyConvex.Analysis.SmoothMBMap.smoothMBOn_iff_smoothMBOn_uncurry}
  \leanok
  \uses{thm:exponential-law, lem:ck-into-smooth, cor:smooth-smooth2, prop:ckk-ck}
  For $f \in \Ck{\infty}(U \times V, F)$, $f^\vee \colon U \to \Ck{\infty}(V, F)$ is smooth, and
  $\Phi \colon \Ck{\infty}(U \times V, F) \to \Ck{\infty}(U, \Ck{\infty}(V, F))$,
  $f \mapsto f^\vee$, is linear and a topological embedding. If $E_2$ is locally compact, $\Phi$
  is an isomorphism of topological vector spaces, and a map $g \colon U \to \Ck{\infty}(V, F)$ is
  smooth if and only if $(x, y) \mapsto g(x)(y)$ is smooth.
\end{theorem}

\begin{proof}
  \leanok
  Smoothness of $f^\vee$ follows from Theorem~\ref{thm:exponential-law} and
  Lemma~\ref{lem:ck-into-smooth}. The embedding is proved as in
  Theorem~\ref{thm:exponential-law}, using that the topology of $\Ck{\infty}(U \times V, F)$ is
  the $\Ck{\infty,\infty}$ topology (Corollary~\ref{cor:smooth-smooth2}).
\end{proof}
